% Part: lambda-calculus
% Chapter: lambda-definability
% Section: partial-recursive-functions

\documentclass[../../../include/open-logic-section]{subfiles}

\begin{document}

\olfileid{lam}{ldf}{par}
\olsection{توابع بازگشتی جزئی \usetoken{S}{lambda definable} هستند}

توابع بازگشتی جزئی توابعی هستند که از توابع پایه با ترکیب، بازگشت
اولیه و کمینه‌سازی نامحدود به دست می‌آیند. تفاوت آن‌ها با توابع
بازگشتی عمومی در این است که لازم نیست توابع به‌کاررفته در جست‌وجوی
نامحدود منظم باشند. حذف شرط منظم‌بودن به این معناست که توابع تعریف‌شده
با کمینه‌سازی ممکن است گاهی تعریف‌نشده باشند.

در نگاه نخست شاید به نظر برسد همان روش‌هایی که نشان می‌دهند همهٔ توابع
بازگشتی عمومیِ تام~!!{lambda definable} هستند، می‌توانند ثابت کنند همهٔ
توابع بازگشتی جزئی نیز~!!{lambda definable} هستند. برای نمونه، اگر
$f$ و~$g$ به‌ترتیب به‌وسیلهٔ ترم‌های~$F$ و~$G$ !!{lambda defined}
باشند، ترکیب~$f$ با~$g$ به‌وسیلهٔ~$\lambd[x][F (G x)]$
!!{lambda defined} است. اما وقتی توابع جزئی باشند، مشکلی پدید می‌آید.
اگر~$g(x)$ تعریف‌نشده باشد، $G x$ صورت نرمالی ندارد. در بیشتر موارد،
این امر به آن معناست که~$F (G x)$ نیز صورت نرمالی ندارد، و این همان
چیزی است که می‌خواهیم. ولی حالتی را در نظر بگیرید که
$F$ برابر~$\lambd[x][\lambd[y][y]]$ است؛ در این صورت~$F (G x)$ صورت
نرمال~$\lambd[y][y]$ را دارد.

این مشکل حل‌ناپذیر نیست و راه‌هایی وجود دارد که همهٔ توابع بازگشتی جزئی
را چنان~!!{lambda define} کنند که مقادیر تعریف‌نشده با ترم‌های فاقد
صورت نرمال نمایش داده شوند. بااین‌حال، این روش‌ها تا اندازه‌ای پیچیده‌تر
و کم‌شهودتر از رویکردی‌اند که برای توابع بازگشتی عمومی در پیش گرفتیم.
قضیه را در اینجا بدون برهان ثبت می‌کنیم:

\begin{thm}
  همهٔ توابع بازگشتی جزئی~!!{lambda definable} هستند.
\end{thm}

\end{document}
